\documentclass{exam}
\usepackage{../../shah311}

\hypersetup{colorlinks=true, linktoc=section, linkcolor=blue}

% p.120: 4.2.2, 4.2.4, 4.2.5, 4.2.7, 4.2.8, 4.2.9, 4.2.10,
% p.126: 4.3.1, 4.3.4, 4.3.9.

\pagestyle{headandfoot}
\firstpageheadrule
\runningheadrule
\firstpageheader{Prof. Rong \\ Real Analysis}{Homework 8}{Jeevan Shah}
\runningheader{Real Analysis \\ Homework 8}{}{Shah}
\firstpagefooter{}{}{}
\runningfooter{ }{\thepage}{ }

\printanswers

\begin{document}
\colorbox{red}{\underline{4.2.2:}}
\begin{solution}
    \begin{parts}
        \part $\displaystyle|f(x) - 9| = |5x - 6 - 9| = |5x - 15| = 3|x - 3| < \epsilon \Rightarrow |x - 3| < \frac{\epsilon}{5} = \delta$.

        \part \begin{align*} 
            |f(x) - 2| = |\sqrt{x} - 2| < 1 &\Rightarrow -1 < \sqrt{x} - 2 < 1 \\
            &\Rightarrow 1 < \sqrt{x} < 3 \\
            &\Rightarrow 1 < x < 9 
        \end{align*}
        Thus we need $|x - 4| < \delta \Rightarrow 1 < x < 9$. Since $|1 - 4| = 3$ and $|9 - 4| = 5$ and $3 < 5$, $\delta = 3$ since if $\delta = 5$ then the interval would be $(-1, 9)$ which is larger than $(1, 9)$.

        \part \begin{align*}
            |f(x) - 3| = |\lfloor x \rfloor - 3| < 1 &\Rightarrow -1 < \lfloor x \rfloor - 3 < 1 \\
            &\Rightarrow 2 < \lfloor x \rfloor < 4 \\
            &\Rightarrow 3 \leq x < 4, \quad \text{since} \lfloor x \rfloor \in \ZZ
        \end{align*}
        Thus, we need $|x - \pi| < \delta \Rightarrow 3 \leq x < 4$. By a similar process as above, ti is clear that $\delta = |3 - \pi|$ for $\epsilon = 1$.

        \part \begin{align*}
            |f(x) - 3| = |\lfloor x \rfloor - 3| < 0.01 &\Rightarrow -0.01 < \lfloor x \rfloor - 3 < 0.01 \\
            &\Rightarrow 2.99 < \lfloor x \rfloor < 3.01 \\
            &\Rightarrow 3 \leq x < 4
        \end{align*}
        So, by the same argument as above, $\delta = |3 - \pi|$.
    \end{parts}
\end{solution}

\colorbox{red}{\underline{4.2.4:}}
\begin{solution}
    \begin{parts}
        \part \begin{align*}
            \left|f(x) - \frac{1}{10}\right| = \left|\frac{1}{\lfloor x \rfloor}- \frac{1}{10}\right| < \frac{1}{2} &\Rightarrow -\frac{1}{2} < \frac{1}{\lfloor x \rfloor} - \frac{1}{10} < \frac{1}{2} \\
            &\Rightarrow -\frac{2}{5} < \frac{1}{\lfloor x \rfloor} < \frac{3}{5} \\
            &\Rightarrow \lfloor x \rfloor > -\frac{5}{2} \quad\text{and}\quad \lfloor x \rfloor > \frac{5}{3} \\
            &\Rightarrow x \geq -2 \quad\text{and}\quad x \geq 2
        \end{align*}
        For an interval centered at $10$, $x \geq 2$ is the tighter bound, so $\delta = 8$.

        \part \begin{align*}
            \left|f(x) - \frac{1}{10}\right| = \left|\frac{1}{\lfloor x \rfloor} - \frac{1}{10}\right| < \frac{1}{50} &\Rightarrow -\frac{1}{50} < \frac{1}{\lfloor x \rfloor} - \frac{1}{10} < \frac{1}{50} \\
            &\Rightarrow \frac{2}{25} < \frac{1}{\lfloor x \rfloor} < \frac{3}{25} \\
            &\Rightarrow \frac{25}{3} < {\lfloor x \rfloor} < \frac{25}{2} \\
            &\Rightarrow 9 \leq x < 13
        \end{align*}
        For an interval centered at $x = 10$, $x \geq 9$ is the tighter bound so $\delta = 1$.

        \part For $9 \leq x < 10, f(x) = \frac{1}{9}$ so $\left|\frac{1}{9} - \frac{1}{10}\right| = \frac{1}{90} = \epsilon$ is the largest values such that no $\delta$ exists. 
    \end{parts}    
\end{solution}

\colorbox{red}{\underline{4.2.5:}}
\begin{solution}
    \begin{parts}
        \part \begin{proof}
            Consider
            \[
                |3x + 4 - 10| = |3x-6| = 3|x-2| < \epsilon \Rightarrow |x-2| < \frac{\epsilon}{3} = \delta. 
            \]
            Choose $\delta = \frac{\epsilon}{3}$. Then, $|x - 2| < \epsilon$ implies 
            \[
                |3x + 4 - 10| < 3\left(\frac{\epsilon}{3}\right) = \epsilon 
            \]
        \end{proof}

        \part \begin{proof}
            Consider 
            \[
                |x^3| < \epsilon \Rightarrow x < \sqrt[3]{\epsilon} = \delta.
            \]
            Choose $\delta = \sqrt[3]{\epsilon}$. Then, by the above $x < \delta \Rightarrow x^3 < \epsilon$.
        \end{proof}

        \part \begin{proof}
            Consider 
            \[
                |x^2 + x - 1 - 5| = |x^2 + x - 6| = |x + 3| \cdot |x - 2| < \epsilon. 
            \]
            Take $\delta \leq 1$ so that $x \in (1, 3)$ and $|x + 3| < 6$. Choose $\delta = \min\left(\frac{\epsilon}{6}, 1\right)$ so that 
            \[
                |x - 2| < \delta \Rightarrow |x+3| \cdot |x-2| < 6\left(\frac{\epsilon}{6}\right) = \epsilon.
            \]
        \end{proof}

        \part \begin{proof}
            Consider 
            \[
                \left|\frac{1}{x} - \frac{1}{3}\right| = \left|\frac{3-x}{3x}\right| < \epsilon.
            \]
            Take $\delta \leq 1$ so $x \in (2, 5)$ and 
            \[
                \left|\frac{3-x}{3x}\right| < \left|\frac{3-x}{6}\right| < \epsilon \Rightarrow |3 - x| < 6\epsilon = \delta.
            \]
            Choose $\delta = \min\left(1, 6\epsilon\right)$ so that 
            \[
                |x - 3| < \delta \Rightarrow \left|\frac{3-x}{3x}\right| < 6\left(\frac{\epsilon}{6}\right) = \epsilon 
            \]
        \end{proof}
    \end{parts}    
\end{solution}

\colorbox{red}{\underline{4.2.7:}}
\begin{solution}
    \begin{proof}
        Assume $\lim_{x\to\,c} g(x) = 0$. That is, $\forall \epsilon > 0, \exists \delta > 0 \st |x - c| < \delta \Rightarrow g(x) < \epsilon$. It follows that $g(x) < \frac{\epsilon}{M}$. Then, $|f(x)g(x)| < \frac{\epsilon}{M}M < \epsilon$ when $|x - c| < \delta$. So, $\lim_{x\to\,c} f(x)g(x) = 0$. 
    \end{proof}    
\end{solution}

\colorbox{red}{\underline{4.2.8:}}
\begin{solution}
    \begin{parts}
        \part We claim the limit does not exist. Consider $x_n = 2 - \frac{1}{n}$ and $y_n = 2 + \frac{1}{n}$ such that $x_n, y_n \to 2$. But, $x_n < 2$ for all $n$ so $f(x_n) \to -1$ since the numerator and denominator will have different signs. Also since $y_n > 2$ for all $n$, $f(y_n) \to 1$. Thus, since there exist two sequences that converge to $2$ but the function values approach different values, the limit cannot exist.

        \part We claim the limit is $-1$. Since $\frac{7}{4} < 2$, the numerator $|x - 2|$ simplifies to $-(x-2)$ in the neighborhood. Thus, for any $\epsilon > 0$ we take $\delta = \frac{1}{4}$ so 
        \[
            0 < \left|x - \frac{7}{4}\right| < \delta \Rightarrow \frac{3}{2} < x < 2
        \]
        and
        \[
            x < 2 \Rightarrow \left|\frac{|x-2|}{x-2} - (-1)\right| = 0
        \]
        as was to be shown.

        \part We claim the limit does not exist. Consider $x_n = \frac{1}{2n}$ and $y_n = \frac{1}{2n+1}$ so $x_n, y_n \to 0$ and $\frac{1}{\lfloor x_n \rfloor} = 2n$ and $\frac{1}{\lfloor y_n \rfloor } = 2n+1$. Then, $(-1)^{1/\lfloor x_n \rfloor} = (-1)^{2n} = 1$ and $(-1)^{\lfloor y_n \rfloor} = (-1)^{2n+1} = -1$. Thus, since there exist two sequences that converge to the same value but have different functional limits, the limit does not exist.

        \part We claim the limit is $0$. Consider $|\sqrt[3]{x}(-1)^{\lfloor 1/x \rfloor} = |\sqrt[3]{x}|$ since $\lim_{x\to\,0} \sqrt[3]{x} = 0$, we have 
        \[
            \lim_{x\to\,0}\sqrt[3]{x}\left(-1\right)^{\lfloor 1/x \rfloor} = 0
        \]
        by the algebraic limit theorem.
    \end{parts}
\end{solution}

\colorbox{red}{\underline{4.2.9:}}
\begin{solution}
    \begin{parts}
        \part Consider $\delta = \frac{1}{M}$. Then $0 < |x - c| < \delta \Rightarrow f(x) = \frac{1}{(1/M)^2} = M^2 > M$.

        \part We define 
        \[
            \lim_{x\to\infty} f(x) = L 
        \]
        if for all $\epsilon > 0$ there exists $N > 0$ such that $x > N$ implies $|f(x) - L| < \epsilon$.
        \begin{proof}
            Take $N = \frac{1}{\epsilon}$. Then $x > N = \frac{1}{\epsilon} \Rightarrow \frac{1}{x} < \epsilon$ since $x > N > 0$ so $\lim_{x \to \infty} \frac{1}{x} = 0$.
        \end{proof}

        \part For every $M > 0$ there exists $N > 0$ such that $x > N$ implies $f(x) > M$. For example, $f(x) = x$ with $N = M$. Then $x > N \Rightarrow f(x) = x > N = M$. Thus $\lim_{x\to\,\infty} f(x) = \infty$.
    \end{parts}    
\end{solution}

\colorbox{red}{\underline{4.2.10:}}
\begin{solution}
    \begin{parts}
        \part Define $\lim_{x \to a^{+}} f(x) = L$ as for all $\epsilon > 0$ there exists $\delta > 0$ such that $0 < x - a < \delta$ implies $|f(x) - L| < \epsilon$ and $\lim_{x \to a^{-}} f(x) = L$ as for all $\epsilon > 0$ there exists $\delta > 0$ such that $0 < a - x < \delta$ implies $|f(x) - L| < \epsilon$. 

        \part \begin{proof}
            For contradiction assume $\lim_{x\to\,a^{+}} f(x) = L$ and $\lim_{x\to\,a^{-}} f(x) = K$ but $\lim_{x\to\,a} f(x) = L$. By (a corresponding definition for sided limits) there exists $x_n$ such that $x_n \to a$ with $x_n < a$ and $f(x_n) \to L$. There also exists $y_n \to a$ with $y_n > a$ and $f(y_n) \to K$. This contradicts the fact that $\lim_{x\to\,a} f(x) = L$ and so both the sided limits must be equal. 
        \end{proof}
    \end{parts}
\end{solution}

\colorbox{red}{\underline{4.3.1:}}
\begin{solution}
    \begin{parts}
        \part Take $\delta = \sqrt[3]{\epsilon}$. Then $x < \delta \Rightarrow x^3 < \delta^3 = \epsilon$.

        \part \begin{proof}
            By the difference of cubes 
            \[
                \sqrt[3]{x} - \sqrt[3]{x} = \frac{x - c}{\sqrt[3]{x^2} - \sqrt[3]{xc} + \sqrt[3]{c^2}}. 
            \]
            Take $\delta_1 < \frac{|c|}{2}$ so $|x - c| < \delta_1$, thus $|x| > |c| - |x-c| > |c| - \frac{|c|}{2} = \frac{|c|}{2}$. Therefore, $|\sqrt[3]{x} \geq \sqrt[3]{\frac{|c|}{2}}$. Then, whenever $|x - c| < \delta_1$ 
            \[
                |\sqrt[3]{x} - \sqrt[3]{c}| \leq \frac{|x-c}{3/4|\sqrt[3]{c}|^2} = \frac{4}{3|c^{2/3}|}|x - c|.
            \]
            Now take 
            \[
                \delta = \min\left(\frac{|c|}{2}, \frac{3}{4}|c|^{2/3}\epsilon\right).
            \]
            So, $|x - c| < \delta$ implies 
            \[
                |\sqrt[3]{x} - \sqrt[3]{c}| \leq \frac{4}{3|c|^{2/3}}|x - c| < \frac{4}{3|c|^{2/3}} \cdot \frac{3}{4}|c|^{2/3}\epsilon = \epsilon
            \]
            and thus $g(x) = \sqrt[3]{x}$ is continuous at some $c \neq 0$. 
        \end{proof}
    \end{parts}    
\end{solution}

\colorbox{red}{\underline{4.3.4:}}
\begin{solution}
    \begin{parts}
        \part Take $f(x) = 0$ and 
        \[
            g(y) = \begin{cases}
                1, &\quad y = 0 \\
                0, &\quad y \neq 0
            \end{cases} 
        \]
        With $p = q = 0$ 
        \[
            \lim_{x \to 0} f(x) = 0 \quad\text{and}\quad \lim_{y \to 0} g(y) = 0
        \]
        Then $g(f(x)) = g(0) = 1$ for all $x$, so 
        \[
            \lim_{x\to\,0} g(f(x)) = 1 \neq 0 = r
        \]

        \part If $f$ and $g$ are continuous with $f(p) = q$ and $g(q) = r$. Since function compositions preserve continuity, 
        \[
            \lim_{x\to\,p} g(f(x)) = g(f(p)) = g(q) = r
        \]

        \part From $(a)$ it is clear that only $f$ continuous is not a strong enough condition. However, only $g$ continuous is since if 
        \[
            \lim_{x\to\,p} f(x) = q \Rightarrow \lim_{x\to\,p} g(f(x)) = g\left(\lim_{x\to\,p} f(x)\right) = g(q) = r
        \]
        since $g$ is continuous at $r$.
    \end{parts}
\end{solution}

\colorbox{red}{\underline{4.3.9}}
\begin{solution}
    It suffices to show $K^c$ is open for $x_0 \in K^c$, $H(x_0) \neq 0$ so we define $\epsilon = \frac{|h(x_0)}{2} > 0$. By definition of continuity, there exists $\delta > 0$ such that $|x - x_0| < \delta$ implies $|h(x) - h(x_0)| < \epsilon$. THus, $|x - x_0| < \delta$ means
    \[
        |h(x)| \geq |h(x_0)| - |h(x) - h(x_0)| > |h(x_0)| - \frac{|h(x_0)|}{2} = \frac{|h(x_0)|}{2} > 0
    \]
    so $h(x) \neq 0 \Rightarrow x \in K^c$ and $|x - x_0| < \delta \subseteq K^c$ so there is an open interval around $x_0$ in $K^c$. It follows that $K^c$ is open and $K$ is then closed. 
\end{solution}
\end{document}